\section*{第 \ref{ch:b3:trace-element-modelling} 章——微量元素模拟}

\begin{solution}{exo:b3:trace-element-modelling:1}
尖晶石二辉橄榄岩：$0.12\times0.40 + 0.88\times0.016 = 0.048 + 0.0141 =
0.062$。石榴子石二辉橄榄岩：$0.195\times0.40 + 0.167\times0.11 +
0.639\times0.016 = 0.078 + 0.0184 + 0.0102 = 0.107$。
\end{solution}

\begin{solution}{exo:b3:trace-element-modelling:2}
$1/(0.05 + 0.02\times0.95) = 1/0.069 = \qty{14.5}{ppm}$；
$1/(0.05 + 0.10\times0.95) = 1/0.145 = \qty{6.9}{ppm}$。
\end{solution}

\begin{solution}{exo:b3:trace-element-modelling:3}
富集程度在 $F \to 0$ 时最大：$1/D = 100$。当 $D = 0$ 时，
$C_L/C_0 = 1/F = 20$。
\end{solution}

\begin{solution}{exo:b3:trace-element-modelling:4}
$\mathrm{La/Yb} = 30/2.5 = 12.0$；$(\mathrm{La}/\mathrm{Yb})_N =
(30/0.232)/(2.5/0.163) = 129.3/15.34 = 8.4$。
\end{solution}

\begin{solution}{exo:b3:trace-element-modelling:5}
\omterm{thm:b3:trace-element-modelling:batch}{批式熔融}：$C_S/C_0 = 0.01/(0.01 + 0.05\times0.99) = 0.168$。\omterm{thm:b3:trace-element-modelling:fractional}{分离熔融}：
$(0.95)^{99} = 0.0062$，低 27 倍。瞬时熔体含
$0.0062/0.01 = 0.62\,C_0$，低于源区。
\end{solution}

\begin{solution}{exo:b3:trace-element-modelling:6}
$D_0 = 0.58\times0.0003 + 0.27\times0.002 + 0.12\times0.053 + 0.03\times0.01
= 0.00737$；$P = 0.20\times0.0003 + 0.25\times0.002 + 0.45\times0.053 +
0.10\times0.01 = 0.0254$。当 $F = 0.02$ 时，\omterm{def:b3:trace-element-modelling:nonmodal}{非实比熔融}为 $1/(0.00737 +
0.02\times0.9746) = 37.2$，实比熔融为 $1/(0.00737 + 0.02\times0.9926) = 36.7$。当
$F = 0.10$ 时，二者分别为 9.54 和 9.38，相差 \qty{1.7}{\percent}。
\end{solution}

\begin{solution}{exo:b3:trace-element-modelling:7}
$D = 0.1$：\omterm{thm:b3:trace-element-modelling:batch}{批式熔融} $1/(0.1 + 0.09) = 5.26$；\omterm{def:b3:trace-element-modelling:fractional}{聚集分离熔体} $(1 - 0.9^{10})/0.1 =
6.51$。$D = 0.01$：\omterm{thm:b3:trace-element-modelling:batch}{批式熔融} $1/(0.01 + 0.099) = 9.17$；\omterm{def:b3:trace-element-modelling:fractional}{聚集分离熔体}
$(1 - 0.9^{100})/0.1 = 10.00$。$D = 0.1$ 时两者差别较大；
而 $D = 0.01$ 时，两种熔体几乎都已含有全部该元素，且两者均趋于
$1/F$。
\end{solution}

\begin{solution}{exo:b3:trace-element-modelling:8}
$D' = 0.99\times0.001 + 0.01 = 0.01099$。首次熔体：$1/D' = 91$，与\omterm{thm:b3:trace-element-modelling:fractional}{分离熔融}的 1000 相比。
当 $F = 0.01$ 时：$91\times0.99^{90.0} = 36.8$，
相比之下 $1000\times0.99^{999} = 0.044$：孔隙熔体使残余体
免于被剥离。
\end{solution}

\begin{solution}{exo:b3:trace-element-modelling:9}
该比值增大的倍数为 $(D_{\mathrm{Sm}} + F(1 - D_{\mathrm{Sm}}))/(D_{\mathrm{La}}
+ F(1 - D_{\mathrm{La}}))$：当 $F = 0.01$ 时，$0.0714/0.0177 = 4.0$；当 $F = 0.10$ 时，
$0.156/0.107 = 1.46$。该比值在 $F$ 较小时对 $F$ 敏感，
在约 \qty{10}{\percent} 以上几乎不敏感；它还依赖于 $D_{\mathrm{Sm}}$，
因而取决于残余矿物。
\end{solution}

\begin{solution}{exo:b3:trace-element-modelling:10}
$\frac1F\int_0^F \frac1D(1 - x)^{1/D - 1}\,\mathrm dx = \frac1F\bigl[-(1 -
x)^{1/D}\bigr]_0^F = [1 - (1 - F)^{1/D}]/F$，即本章给出的公式；若
$(1 - F)^{1/D} \to 0$，则其趋于 $1/F$ 乘以 $C_0$。在 \qty{1}{\percent} 以内：
$(1 - F)^{100} < 0.01$，$F > 1 - 0.01^{0.01} = 0.045$，约为
\qty{4.5}{\percent}。
\end{solution}

\begin{solution}{exo:b3:trace-element-modelling:11}
斜率 $a = (10 - 6)/(20 - 10) = \qty{0.40}{ppm^{-1}}$；截距 $b = 10 -
0.40\times20 = 2.0$。$D_{\mathrm{Yb}} = 0.40\times0.44 = 0.176$；
$C_0^{\mathrm{La}} = 2.0\times0.44/0.824 = \qty{1.07}{ppm}$。$F_P = 1.07/20 =
\qty{5.3}{\percent}$，$F_Q = 1.07/10 = \qty{10.7}{\percent}$。
\end{solution}

\begin{solution}{exo:b3:trace-element-modelling:12}
$F = x_i/p_i$：单斜辉石 $0.12/0.45 = 0.27$，尖晶石 $0.03/0.10 = 0.30$，
斜方辉石 $0.27/0.25 = 1.08$，橄榄石 $0.58/0.20 = 2.9$；单斜辉石
最先熔融。当 $F = 0.2$ 时：$\bar D = (0.00737 - 0.0254\times0.2)/0.8 =
0.0029$。超过 $F = 0.27$ 后，$p_i$ 发生变化（已无单斜辉石可熔），
而假定 $p_i$ 为常数的公式不再适用。
\end{solution}

\begin{solution}{pb:b3:trace-element-modelling:1}
\textbf{1.} $D_{\mathrm{La}} = 0.12\times0.058 + 0.88\times0.001 = 0.00784$；
$D_{\mathrm{Yb}} = 0.12\times0.58 + 0.88\times0.16 = 0.210$。

\textbf{2.} $D_{\mathrm{La}} = 0.195\times0.058 + 0.167\times0.001 +
0.639\times0.001 = 0.0121$；$D_{\mathrm{Yb}} = 0.195\times0.58 +
0.167\times4.5 + 0.639\times0.16 = 0.967$。取整后为 0.0078、0.21、0.012、0.97。

\textbf{3.} $D_{\mathrm{Yb}}$ 由 0.21 升至 0.97，因为石榴子石富集
重稀土（$D = 4.5$）。镱仍是极弱的不相容元素，
$0.97 < 1$。

\textbf{4.} $(0.648/0.232)/(0.441/0.163) = 2.79/2.71 = 1.03$：平坦型模式，
约为球粒陨石的 2.8 倍。

\textbf{5.} $0.648/(0.0078 + 0.01\times0.9922) = \qty{36.6}{ppm}$；在
\qty{5}{\percent} 时为 \qty{11.3}{ppm}；在 \qty{10}{\percent} 时为 \qty{6.05}{ppm}。

\textbf{6.} $F = (0.648/9.6 - 0.0078)/0.9922 = (0.0675 - 0.0078)/0.9922 =
0.0602$，即 \qty{6.0}{\percent}。

\textbf{7.} $\mathrm{Yb} = 0.441/(0.21 + 0.0602\times0.79) = \qty{1.71}{ppm}$，
La/Yb $= 5.6$；熔岩 A 的值为 1.7 和 5.6：与尖晶石源区相符。

\textbf{8.} La $= 0.0078\times9.6 = \qty{0.075}{ppm}$，Yb $= 0.21\times1.71 =
\qty{0.36}{ppm}$。剩余的镧：$0.940\times0.075/0.648 = 0.11$，约为
\qty{11}{\percent}。

\textbf{9.} $F = 0.648/9.6 = \qty{6.75}{\percent}$.

\textbf{10.} $(0.648/0.0078)(1 - F)^{127.2}$：在 \qty{0.1}{\percent} 时为 \qty{73}{ppm}，
在 \qty{1}{\percent} 时为 \qty{23}{ppm}，在
\qty{6}{\percent} 时为 \qty{0.03}{ppm}。

\textbf{11.} $0.648\times(1 - 0.940^{128.2})/0.060 = \qty{10.8}{ppm}$.

\textbf{12.} 取 $F = \qty{6.75}{\percent}$，在此 $F$ 下，$(1 - F)^{1/D} =
\num{1.3e-4}$，因此汇集熔体含有源区的全部镧，并且
$\bar C_L = C_0/F$，与完全不相容元素的情形相同。

\textbf{13.} $0.648\times0.9325^{127.2} = \qty{9e-5}{ppm}$，比\omterm{thm:b3:trace-element-modelling:batch}{批式熔融}的
\qty{0.075}{ppm} 约低 800 倍。

\textbf{14.} 两种模型给出答案的范围：\omterm{thm:b3:trace-element-modelling:batch}{批式熔融}为 \qty{6.0}{\percent}，
汇集\omterm{thm:b3:trace-element-modelling:fractional}{分离熔融}为 \qty{6.75}{\percent}。镧可将 $F$ 确定到
约 \qty{1}{\percent} 的精度，与所用模型无关。

\textbf{15.} 熔岩 A：La/Yb $= 5.6$，$(\mathrm{La}/\mathrm{Yb})_N =
(9.6/0.232)/(1.7/0.163) = 4.0$。熔岩 B：La/Yb $= 34.7$，$(\mathrm{La}/\mathrm{Yb})_N
= (15.6/0.232)/(0.45/0.163) = 24.4$。

\textbf{16.} 当 $F \to 0$ 时：$(0.648/0.0078)/(0.441/0.21) = 83.1/2.10 = 39.6$。

\textbf{17.} 此时 $F = \qty{0.11}{\percent}$，熔体将含
\qty{72}{ppm} 的 La 和 \qty{2.1}{ppm} 的 Yb。熔岩 B 含 \qty{15.6}{ppm} 的 La：
尖晶石源区无法生成它。

\textbf{18.} $F = (0.648/15.6 - 0.012)/0.988 = 0.0299$，即 \qty{3.0}{\percent}；
$\mathrm{Yb} = 0.441/(0.97 + 0.0299\times0.03) = \qty{0.454}{ppm}$，与
实测值相符。

\textbf{19.} 在约 60 至 \qty{90}{km} 处，石榴子石取代尖晶石：熔岩 B 的熔融始于
更深处，熔岩 A 则较浅，位于尖晶石区。石榴子石二辉橄榄岩熔体若含
\qty{9.6}{ppm} 的 La（$F \approx \qty{5.6}{\percent}$），其 La/Yb
应接近 21，而非 5.6：熔岩 A 并非来自石榴子石二辉橄榄岩。

\textbf{20.} $a = 0.97/0.441 = \qty{2.20}{ppm^{-1}}$，$b = 0.648\times0.03/0.441
= 0.044$；当 La $= 15.6$ 时，La/Yb $= 2.20\times15.6 + 0.044 = 34.4$，与
熔岩 B 的 34.7 相比。

\textbf{21.} 橄榄石对两种元素都不接受：La/Yb 保持不变。原始
熔体含 La $9.6\times0.9 = \qty{8.64}{ppm}$，因此 $F = (0.648/8.64 -
0.0078)/0.9922 = \qty{6.8}{\percent}$，而非 \qty{6.0}{\percent}。

\textbf{22.} 同一源区可以生成两种熔岩：B 是较小的熔体
（\qty{3}{\percent}），始于石榴子石区，深于约 60 至
\qty{90}{km}；A 则是位于尖晶石区的较大熔体。

\textbf{23.} 熔岩 A 是尖晶石二辉橄榄岩的 \textbf{\qty[text-series-to-math=true]{6.0}{\percent}} \omterm{thm:b3:trace-element-modelling:batch}{批式熔融}产物
（若熔体是经分离抽取后汇集而成，则为
\qty{6.75}{\percent}）。
\end{solution}
